% Part: first-order-logic
% Chapter: axiomatic-deduction
% Section: proving-things

\documentclass[../../../include/open-logic-section]{subfiles}

\begin{document}

\iftag{FOL}
      {\olfileid{fol}{axd}{pro}}
      {\olfileid{pl}{axd}{pro}}

\olsection{أمثلة على \usetoken{P}{derivation}}


\begin{ex}
افترض أننا نريد إثبات $(\lnot !D \lor !E) \lif (!D \lif
!E)$. من الواضح أن هذه ليست حالة من أي بديهية لدينا، ولذلك علينا استخدام
قاعدة \MP{} لكي !!{derive} إياها. والقاعدة الوحيدة التي نحتاج إليها هنا
هي إثبات المقدَّم، وهي تتيح لنا، عند إعطائنا $!A$ و$!A \lif !B$، تبرير~$!B$.
ومن الاستراتيجيات الممكنة استخدام \olref[prp]{ax:lor3} بحيث تكون $!A$ هي
$\lnot !D$، و$!B$ هي $!E$، و$!C$ هي $!D \lif !E$، أي الحالة
\[
(\lnot !D \lif (!D \lif !E)) \lif ((!E \lif (!D \lif !E)) \lif ((\lnot
!D \lor !E) \lif (!D \lif !E))).
\]
لماذا؟ لأن تطبيقين لقاعدة إثبات المقدَّم ينتجان الجزء الأخير، وهو ما نريده. ونرى
بسهولة أن $\lnot !D \lif (!D \lif !E)$ حالة من
\olref[prp]{ax:lnot2}، وأن $!E \lif (!D \lif !E)$ حالة من
\olref[prp]{ax:lif1}. ومن ثم يكون !!{derivation} المطلوب:
\begin{derivation}
  1. &  $\lnot !D \lif (!D \lif !E)$ & \olref[prp]{ax:lnot2} \\
  2. & $(\lnot !D \lif (!D \lif !E)) \lif {}$\\
  &\qquad $((!E \lif (!D \lif !E)) \lif ((\lnot !D \lor !E) \lif (!D \lif !E)))$ & \olref[prp]{ax:lor3}\\
  3. & $(!E \lif (!D \lif !E)) \lif ((\lnot !D \lor !E) \lif (!D \lif !E))$ & 1, 2, \MP\\
  4. & $!E \lif (!D \lif !E)$ & \olref[prp]{ax:lif1}\\
  5. & $(\lnot !D \lor !E) \lif (!D \lif !E)$ & 3, 4, \MP
\end{derivation}
\end{ex}

\begin{ex}\ollabel{ex:identity}
لنحاول إيجاد !!a{derivation} لــ$!D \lif !D$. إنها ليست حالة من بديهية،
ولذلك علينا استخدام \MP{} لكي !!{derive} إياها.
\olref[prp]{ax:lif1} بديهية على الصورة~$!A \lif !B$ يمكننا تطبيق~\MP
عليها. ولكي تكون مفيدة، يجب طبعًا أن تكون $!B$ التي تبررها \MP{} خطوةً
صحيحة في هذه الحالة هي~$!D \lif !D$، لأن هذا هو ما نريد !!{derive}.
وهذا يعني أن $!A$ يجب أن تكون أيضًا $!D$، أي قد ننظر في الحالة الآتية من
\olref[prp]{ax:lif1}:
\[
!D \lif (!D \lif !D)
\]
ولكي نطبق \MP، نحتاج أيضًا إلى تبرير المقدمة الثانية الموافقة، وهي~$!A$.
لكنها ستكون في حالتنا~$!D$، ولن نستطيع !!{derive}~$!D$ وحدها. ولذلك
نحتاج إلى استراتيجية مختلفة.

البديهية الأخرى التي لا تتضمن إلا~$\lif$ هي \olref[prp]{ax:lif2}، أي
\[
(!A \lif (!B \lif !C)) \lif ((!A \lif !B) \lif (!A \lif !C))
\]
يمكننا الوصول إلى الشرطية المتداخلة الأخيرة بتطبيق \MP{} مرتين. وهذا يعني
مرة أخرى أننا نريد حالة من \olref[prp]{ax:lif2} تكون فيها $!A \lif !C$
هي $!D \lif !D$، أي ال!!{formula} التي نهدف إليها. وعندئذ تكون كل من
$!A$ و$!C$ هي~$!D$ طبعًا. فكيف نختار~$!B$ بحيث تكون كل من
$!A \lif (!B \lif !C)$ و$!A \lif !B$، أي في حالتنا
$!D \lif (!B \lif !D)$ و$!D \lif !B$، !!{derivable} أيضًا؟ الصيغة
الأولى حالة من \olref[prp]{ax:lif1} بالفعل أيًا يكن ما نختاره لــ$!B$.
وتكون $!D \lif !B$ حالة أخرى من \olref[prp]{ax:lif1} إذا كانت $!B$ هي
$(!D \lif !D)$. ومن ثم يكون !!{derivation} المطلوب:
\begin{derivation}
  1. & $!D \lif ((!D \lif !D) \lif !D)$ & \olref[prp]{ax:lif1}\\
  2. & $(!D \lif ((!D \lif !D) \lif !D)) \lif {}$\\
  & \qquad $((!D \lif (!D \lif !D)) \lif (!D \lif !D))$ & \olref[prp]{ax:lif2}\\
  3. & $(!D \lif (!D \lif !D)) \lif (!D \lif !D)$ & 1, 2, \MP\\
  4. & $!D \lif (!D \lif !D)$ & \olref[prp]{ax:lif1}\\
  5. & $!D \lif !D$ & 3, 4, \MP
\end{derivation}
\end{ex}

\begin{ex}\ollabel{ex:chain}
نريد أحيانًا أن نبين وجود !!a{derivation} ل!!{formula} ما من
!!{formula}s أخرى~$\Gamma$. فمثلًا، لنبين أنه يمكننا !!{derive}
$!A \lif !C$ من $\Gamma = \{!A \lif !B, !B \lif !C\}$.
\begin{derivation}
  1. & $!A \lif !B$ & \Hyp\\
  2. & $!B \lif !C$ & \Hyp\\
  3. & $(!B \lif !C) \lif (!A \lif (!B \lif !C))$ & \olref[prp]{ax:lif1} \\
  4. & $!A \lif (!B \lif !C)$ & 2, 3, \MP\\
  5. & $(!A \lif (!B \lif !C)) \lif {}$\\
  & \qquad $((!A \lif !B) \lif (!A \lif !C))$ & \olref[prp]{ax:lif2}\\
  6. & $((!A \lif !B) \lif (!A \lif !C))$ & 4, 5, \MP\\
  7. & $!A \lif !C$ & 1, 6, \MP
\end{derivation}
تشير الأسطر الموسومة «\Hyp» (اختصارًا لــ«فرضية») إلى أن ال!!{formula}
في ذلك السطر !!a{element} من~$\Gamma$.
\end{ex}

\begin{prop}
  \ollabel{prop:chain} إذا كان $\Gamma \Proves !A \lif !B$ و$\Gamma
  \Proves !B \lif !C$، فإن $\Gamma \Proves !A \lif !C$.
\end{prop}

\begin{proof}
  افترض أن $\Gamma \Proves !A \lif !B$ و$\Gamma \Proves !B \lif
  !C$. عندئذ يوجد !!a{derivation} لــ$!A \lif !B$ من~$\Gamma$؛ ويوجد
  كذلك !!a{derivation} لــ~$!B \lif !C$ من~$\Gamma$. اجمعهما في
  !!{derivation} واحد بوصل إحداهما بالأخرى. ثم أضف الأسطر 3--7 من
  ال!!{derivation} في المثال السابق. وهذا !!a{derivation} لــ
  $!A \lif !C$---وهي السطر الأخير من ال!!{derivation} الجديد---من
  ~$\Gamma$. لاحظ أن تبريري السطرين 4 و~7 يظلان صحيحين إذا استُبدلت
  بالإحالة إلى السطر رقم~1 إحالة إلى السطر الأخير من ال!!{derivation}
  لــ~$!A \lif !B$، وبالإحالة إلى السطر رقم~2 إحالة إلى السطر الأخير
  من ال!!{derivation} لــ~$!B \lif !C$.
\end{proof}

\begin{prob}
بيّن أن العبارات الآتية صحيحة بعرض !!{derivation}s من البديهيات:
\begin{enumerate}
\item $(!A \land !B) \lif (!B \land !A)$
\item $((!A \land !B) \lif !C) \lif (!A \lif (!B \lif !C))$
\item $\lnot(!A \lor !B) \lif \lnot !A$.
\end{enumerate}
\end{prob}



\end{document}
